\documentclass[11pt,a4paper]{article}
\usepackage[T1]{fontenc}
\usepackage[utf8]{inputenc}
\usepackage[french]{babel}
\usepackage{lmodern}
\usepackage{textcomp}
\usepackage{amsmath,amssymb}
\usepackage{geometry}
\usepackage{tabularx,array}
\usepackage[table]{xcolor}
\usepackage{tikz}
\geometry{left=8mm,right=8mm,top=7mm,bottom=7mm}
\setlength{\parindent}{0pt}
\pagestyle{empty}
\setlength{\tabcolsep}{3pt}
\renewcommand{\arraystretch}{1.05}
\renewcommand{\tabularxcolumn}[1]{m{#1}}

% --- Palette pensée pour une impression noir et blanc ---
\definecolor{canlight}{gray}{0.93}
\definecolor{canmid}{gray}{0.72}
\definecolor{candark}{gray}{0.22}
\definecolor{answerblue}{RGB}{0,82,165}

% --- Colonnes ---
\newcolumntype{Q}{>{\centering\arraybackslash}X}
\newcolumntype{N}{>{\centering\arraybackslash}m{6mm}}
\newcolumntype{R}{>{\centering\arraybackslash}m{22mm}}

% --- En-tête ---
\newcommand{\canheader}[1]{%
  \begin{center}
    {\LARGE\bfseries COURSE AUX NOMBRES}\par
    \vspace{0.7mm}
    {\color{canmid}\rule{27mm}{0.45pt}\hspace{2mm}\color{black}$\diamond$\hspace{1.4mm}{\large\bfseries #1}\hspace{1.4mm}$\diamond$\hspace{2mm}\color{canmid}\rule{27mm}{0.45pt}\color{black}}
  \end{center}
  \vspace{0.8mm}
  \renewcommand{\arraystretch}{1.25}%
  \begin{tabularx}{\linewidth}{>{\centering\arraybackslash}X >{\centering\arraybackslash}X >{\centering\arraybackslash}m{40mm}}
    \rowcolor{canlight}
    \footnotesize\textbf{30 questions} -- 9 minutes &
    \footnotesize\textbf{Sans calculatrice} &
    \footnotesize\textbf{Score :} \rule{13mm}{0.45pt}\,/30
  \end{tabularx}
  \renewcommand{\arraystretch}{1.05}%
  \vspace{1mm}
  \begin{center}
    {\scriptsize Réponse uniquement \quad$\diamond$\quad Reporter le score sur la fiche annuelle.}
  \end{center}
  \vspace{1.2mm}
}

\newcommand{\qrow}[2]{\textbf{#1} & \parbox[c][15.0mm][c]{\linewidth}{\centering #2} & \\ \hline}
\newcommand{\qrowfig}[2]{\textbf{#1} & \parbox[c][27mm][c]{\linewidth}{\centering #2} & \\ \hline}
\newcommand{\qrowcorr}[3]{\textbf{#1} & \parbox[c][15.0mm][c]{\linewidth}{\centering #2} & {\color{answerblue}\bfseries #3} \\ \hline}
\newcommand{\qrowfigcorr}[3]{\textbf{#1} & \parbox[c][27mm][c]{\linewidth}{\centering #2} & {\color{answerblue}\bfseries #3} \\ \hline}

\newcommand{\cantable}[1]{%
\arrayrulecolor{canmid}
\begin{tabularx}{\linewidth}{N|Q|R}
\rowcolor{canlight}\textbf{N\textsuperscript{o}}&\centering\textbf{Question}&\textbf{Réponse}\\\hline
#1
\end{tabularx}
\arrayrulecolor{black}
}

\newcommand{\hypotfigure}{%
\begin{tikzpicture}[baseline=(current bounding box.center),scale=0.72,line width=0.6pt]
  \node[font=\bfseries\normalsize] at (1.6,2.55) {Quelle est l'hypoténuse ?};
  \coordinate (A) at (0,0); \coordinate (B) at (0,1.85); \coordinate (C) at (3.2,0);
  \draw (A)--(B)--(C)--cycle;
  \draw (0,0) rectangle (0.32,0.32);
  \node[below left] at (A) {\small A};
  \node[above left] at (B) {\small B};
  \node[below right] at (C) {\small C};
\end{tikzpicture}}


\newcommand{\canheadercorr}[1]{%
  \begin{center}
    {\LARGE\bfseries COURSE AUX NOMBRES}\par
    \vspace{0.7mm}
    {\color{canmid}\rule{27mm}{0.45pt}\hspace{2mm}\color{black}$\diamond$\hspace{1.4mm}{\large\bfseries #1}\hspace{1.4mm}$\diamond$\hspace{2mm}\color{canmid}\rule{27mm}{0.45pt}\color{black}}
  \end{center}
  \vspace{0.8mm}
  \renewcommand{\arraystretch}{1.25}%
  \begin{tabularx}{\linewidth}{>{\centering\arraybackslash}X >{\centering\arraybackslash}X >{\centering\arraybackslash}m{40mm}}
    \rowcolor{canlight}
    \footnotesize\textbf{30 questions} -- 9 minutes &
    \footnotesize\textbf{Sans calculatrice} &
    \footnotesize\textcolor{answerblue}{\textbf{CORRIGÉ}}
  \end{tabularx}
  \renewcommand{\arraystretch}{1.05}%
  \vspace{1mm}
  \begin{center}
    {\scriptsize Les réponses apparaissent en \textcolor{answerblue}{\textbf{bleu}}.}
  \end{center}
  \vspace{1.2mm}
}


\begin{document}
\canheader{CAN 1}

\begin{minipage}[t]{0.493\textwidth}
\normalsize
\cantable{
\qrow{1}{$7\times 8$}
\qrow{2}{Dans $52{,}374$, quel est le chiffre des centièmes ?}
\qrow{3}{$-4+9$}
\qrow{4}{Compléter : $4\times\ldots=36$}
\qrow{5}{$25\,\%$ de 80}
\qrow{6}{Le tiers du quadruple de 18}
\qrow{7}{$46+29$}
\qrow{8}{Donner sous forme de fraction le nombre qui, multiplié par 5, donne 3.}
\qrow{9}{$4{,}07\times100$}
\qrow{10}{Moyenne de 8 ; 10 ; 12}
\qrow{11}{Le plus grand : $2{,}7$ ou $2{,}07$ ?}
\qrow{12}{L'opposé de $-7$}
\qrow{13}{$\dfrac34$ de 20}
\qrow{14}{$1{,}5$ h = $\ldots$ h $\ldots$ min}
\qrow{15}{$3+7\times2$}
}
\end{minipage}\hfill
\begin{minipage}[t]{0.493\textwidth}
\normalsize
\cantable{
\qrow{16}{Simplifier $\dfrac{18}{24}$}
\qrow{17}{Périmètre d'un rectangle de 7 cm sur 3 cm}
\qrow{18}{Une valeur a une fréquence de $0{,}2$ dans une série de 30 données. Son effectif est :}
\qrow{19}{Écriture décimale de $\dfrac14$}
\qrow{20}{Une aire de $24\,\mathrm{cm}^2$ correspond à combien de carrés de $1\,\mathrm{cm}^2$ ?}
\qrow{21}{486 est-il divisible par 3 ?}
\qrow{22}{$0{,}5\times36$}
\qrow{23}{Un volume s'exprime en : cm, cm$^2$ ou cm$^3$ ?}
\qrow{24}{3 kg de pommes coûtent 6 \texteuro. Combien coûtent 5 kg ?}
\qrowfig{25}{\hypotfigure}
\qrow{26}{Quel est le côté d'un carré d'aire $49\,\mathrm{cm}^2$ ?}
\qrow{27}{$19\times5$ est-il plus proche de 50, de 100 ou de 500 ?}
\qrow{28}{Encadrer $7{,}6$ par deux entiers consécutifs.}
\qrow{29}{Réduire : $3x+2x$}
\qrow{30}{Pour $x=4$, calculer $2x+3$.}
}
\end{minipage}

\clearpage

\canheadercorr{CAN 1}

\begin{minipage}[t]{0.493\textwidth}
\normalsize
\cantable{
\qrowcorr{1}{$7\times 8$}{$56$}
\qrowcorr{2}{Dans $52{,}374$, quel est le chiffre des centièmes ?}{$7$}
\qrowcorr{3}{$-4+9$}{$5$}
\qrowcorr{4}{Compléter : $4\times\ldots=36$}{$9$}
\qrowcorr{5}{$25\,\%$ de 80}{$20$}
\qrowcorr{6}{Le tiers du quadruple de 18}{$24$}
\qrowcorr{7}{$46+29$}{$75$}
\qrowcorr{8}{Donner sous forme de fraction le nombre qui, multiplié par 5, donne 3.}{$\dfrac35$}
\qrowcorr{9}{$4{,}07\times100$}{$407$}
\qrowcorr{10}{Moyenne de 8 ; 10 ; 12}{$10$}
\qrowcorr{11}{Le plus grand : $2{,}7$ ou $2{,}07$ ?}{$2{,}7$}
\qrowcorr{12}{L'opposé de $-7$}{$7$}
\qrowcorr{13}{$\dfrac34$ de 20}{$15$}
\qrowcorr{14}{$1{,}5$ h = $\ldots$ h $\ldots$ min}{1 h 30 min}
\qrowcorr{15}{$3+7\times2$}{$17$}
}
\end{minipage}\hfill
\begin{minipage}[t]{0.493\textwidth}
\normalsize
\cantable{
\qrowcorr{16}{Simplifier $\dfrac{18}{24}$}{$\dfrac34$}
\qrowcorr{17}{Périmètre d'un rectangle de 7 cm sur 3 cm}{20 cm}
\qrowcorr{18}{Une valeur a une fréquence de $0{,}2$ dans une série de 30 données. Son effectif est :}{$6$}
\qrowcorr{19}{Écriture décimale de $\dfrac14$}{$0{,}25$}
\qrowcorr{20}{Une aire de $24\,\mathrm{cm}^2$ correspond à combien de carrés de $1\,\mathrm{cm}^2$ ?}{$24$}
\qrowcorr{21}{486 est-il divisible par 3 ?}{oui}
\qrowcorr{22}{$0{,}5\times36$}{$18$}
\qrowcorr{23}{Un volume s'exprime en : cm, cm$^2$ ou cm$^3$ ?}{cm$^3$}
\qrowcorr{24}{3 kg de pommes coûtent 6 \texteuro. Combien coûtent 5 kg ?}{10 \texteuro}
\qrowfigcorr{25}{\hypotfigure}{[BC]}
\qrowcorr{26}{Quel est le côté d'un carré d'aire $49\,\mathrm{cm}^2$ ?}{7 cm}
\qrowcorr{27}{$19\times5$ est-il plus proche de 50, de 100 ou de 500 ?}{$100$}
\qrowcorr{28}{Encadrer $7{,}6$ par deux entiers consécutifs.}{$7<7{,}6<8$}
\qrowcorr{29}{Réduire : $3x+2x$}{$5x$}
\qrowcorr{30}{Pour $x=4$, calculer $2x+3$.}{$11$}
}
\end{minipage}
\end{document}
